\title{DeepSeekMath: Pushing the Limits of Mathematical Reasoning in Open Language Models}

\begin{abstract}
The task of mathematical reasoning remains particularly demanding for language models, owing to the inherent complexity and structure of mathematics. This work presents DeepSeekMath~7B, an extension of DeepSeek-Coder-Base-v1.5 7B that undergoes further pre-training with 120B tokens pertinent to mathematics, obtained from Common Crawl and supplemented with natural language and programming data. Remarkably, DeepSeekMath~7B attains a score of 51.7\% on the competition-grade MATH benchmark, achieving results that are comparable to those of Gemini-Ultra and GPT-4, all without utilizing external toolchains or voting strategies. Employing self-consistency sampling with 64 outputs, DeepSeekMath~7B reaches 60.9\% on the MATH dataset. This enhancement in mathematical reasoning proficiency is ascribed to two main innovations: (1) a rigorously constructed data filtering mechanism that leverages the breadth of math-related content in web datasets, and (2) the development of Group Relative Policy Optimization (GRPO), a modification of Proximal Policy Optimization (PPO), which both bolsters mathematical reasoning capability and improves the efficiency of PPO's memory consumption.
\end{abstract}

\section{Introduction}

The advent of large language models (LLMs) has transformed methods of mathematical reasoning within artificial intelligence, leading to notable progress in benchmarks for quantitative reasoning~\citep{MATH} as well as for geometric problem solving~\citep{trinh2024solving}. Additionally, such models have become valuable tools in supporting individuals tackling sophisticated mathematical challenges~\citep{tao}. Despite these advances, high-performing systems like GPT-4~\citep{gpt4} and Gemini-Ultra~\citep{gemini} remain inaccessible to the public, while existing open-source models exhibit a marked performance gap relative to their proprietary counterparts.

This paper presents DeepSeekMath, a specialized language model tailored for mathematics that exhibits markedly superior mathematical reasoning abilities compared to other open-source alternatives and demonstrates results approaching those of GPT-4 on academic tasks. Central to this achievement is the creation of the DeepSeekMath Corpus, an extensive, high-quality pre-training resource encompassing 120B math-related tokens. The corpus is curated from the Common Crawl (CC) through the application of a fastText-based classifier \citep{joulin2016fasttext}. Initially, the classifier is built using positive samples drawn from OpenWebMath \citep{openwebmath} and negative samples sourced from a broad array of unrelated web documents. The classifier then facilitates the identification of further math-relevant examples from the CC, with additional manual annotation to ensure higher quality. This augmented collection is subsequently used to refine the classifier, boosting its efficacy. Experimental results affirm the strong quality of the resulting corpus, as our foundational DeepSeekMath-Base 7B model achieves 64.2\% accuracy on GSM8K \citep{gsm8k} and 36.2\% on the advanced MATH dataset \citep{MATH}, thereby surpassing Minerva 540B \citep{minerva}. The multilingual scope of the DeepSeekMath Corpus also translates to noticeable gains on Chinese mathematics benchmarks \citep{wei2023cmath,agieval}. We view our work in mathematical data curation and processing as an initial step for the broader research community, anticipating considerable further advancements in this area.

DeepSeekMath-Base is constructed by initializing with DeepSeek-Coder-Base-v1.5 7B \citep{deepseek-coder}, as empirical evidence shows that models pre-trained on code perform better than those derived from general-purpose LLMs. Additionally, we find that training on mathematical data enhances the model's performance not only in mathematics but also on broader benchmarks such as MMLU \citep{mmlu} and BBH \citep{bbh}, demonstrating a significant boost in general reasoning as well as mathematical proficiency.

Subsequent to pre-training, we implement mathematical instruction tuning on DeepSeekMath-Base, incorporating chain-of-thought \citep{cot}, program-of-thought \citep{pot,pal}, and reasoning tasks requiring tool use \citep{tora}. The outcome is DeepSeekMath-Instruct 7B, which surpasses all other 7B models and achieves performance levels competitive with 70B-sized open-source models that are instruction-tuned.

In addition, we propose Group Relative Policy Optimization (GRPO), an RL algorithm that modifies Proximal Policy Optimization (PPO) \citep{schulman2017proximal}. GRPO eliminates the need for a critic model, instead inferring baselines from group-level scores, thereby dramatically decreasing computational requirements. Using only a portion of the English instruction tuning set, GRPO delivers notable gains over DeepSeekMath-Instruct, with substantial improvements recorded on both in-domain (e.g., GSM8K: 82.9\% $\rightarrow$ 88.2\%, MATH: 46.8\% $\rightarrow$ 51.7\%) and out-of-domain mathematics tasks (for instance, CMATH: 84.6\% $\rightarrow$ 88.8\%) after the RL phase. We also introduce a unified framework that situates various techniques, including Rejection Sampling Fine-Tuning (RFT) \citep{yuan2023scaling}, Direct Preference Optimization (DPO) \citep{dpo}, PPO, and GRPO, showing that these approaches can be interpreted as forms of direct or streamlined RL. Comprehensive experiments are conducted—exploring facets such as online versus offline optimization, outcome-based versus process-based supervision, and single-turn compared to iterative RL—to thoroughly analyze key properties of the paradigm. Finally, we offer an explanation of the mechanisms by which our RL strategies enhance instruction-tuned model performance, and we highlight prospective pathways for realizing more powerful RL via this unified conceptual foundation.

\subsection{Contributions}
This work advances large-scale mathematical pre-training and delves into the application and assessment of reinforcement learning.

\textbf{Math Pre-Training at Scale}

\begin{itemize}[topsep=0pt]
    \item We demonstrate that the openly available Common Crawl corpus possesses significant mathematical information. Through a carefully engineered data curation pipeline, we assembled the DeepSeekMath Corpus: a high-quality collection comprising 120B tokens extracted from web content screened for mathematical relevance. This corpus substantially surpasses the math-focused web datasets used in Minerva \citep{minerva} by nearly sevenfold and exceeds the size of OpenWebMath \citep{openwebmath} by a factor of nine.

    \item Our foundational model, DeepSeekMath-Base 7B, performs on par with Minerva 540B \citep{minerva}, illustrating that parameter count alone does not determine success in mathematical reasoning tasks. In fact, a more compact model, when trained on superior data, can also exhibit robust performance.

    \item Our empirical observations from math pre-training experiments reveal that prior training on code enhances a model’s mathematical problem-solving abilities, regardless of tool usage. This finding partially addresses the long-debated issue: \textit{does code pre-training augment reasoning skills?} Our results affirm this, at least concerning mathematical reasoning.

    \item Despite the prevalent use of arXiv paper data for math-related model training, our evaluation indicates that it confers no measurable benefits across all the mathematical benchmarks examined in this study.
\end{itemize}

\textbf{Exploration and Analysis of Reinforcement Learning}

\begin{itemize}[topsep=0pt]
    \item We present Group Relative Policy Optimization (GRPO), a novel reinforcement learning framework that achieves high efficiency and effectiveness. Unlike Proximal Policy Optimization (PPO), GRPO omits the need for a critic network by leveraging group-based score estimates to approximate the baseline, substantially cutting down the computational cost during training.
    \item Our empirical results show that employing GRPO considerably boosts the outcomes of our instruction-tuned model, DeepSeekMath-Instruct, utilizing only instruction-tuning datasets. In addition, we find that reinforcement learning with GRPO leads to noticeable improvements in generalization to out-of-domain tasks.
    \item We propose a comprehensive framework to relate various algorithms, including RFT, DPO, PPO, and GRPO. Through systematic experimentation—comparing, for instance, online versus offline training paradigms, supervision by outcomes versus processes, and single-turn versus iterative reinforcement learning—we analyze and elucidate the critical factors underpinning this unified approach.
    \item Grounded in this integrative paradigm, we delve into the underlying causes driving reinforcement learning efficacy, and outline several promising pathways for enhancing reinforcement learning in large language models.
\end{itemize}

\subsection{Summary of Evaluations and Metrics}

\begin{itemize}[topsep=0pt]
    \item \textbf{English and Chinese Mathematical Reasoning}: Our evaluation protocol encompasses thorough testing on both English and Chinese benchmarks, spanning mathematical challenges appropriate for learners from grade school up through the collegiate level. English assessments draw upon datasets such as GSM8K \citep{gsm8k}, MATH \citep{MATH}, SAT \citep{llemma}, OCW Courses \citep{minerva}, and MMLU-STEM \citep{mmlu}. For Chinese, benchmarks include MGSM-zh \citep{mgsm}, CMATH \citep{wei2023cmath}, Gaokao-MathCloze \citep{agieval}, and Gaokao-MathQA \citep{agieval}. We systematically measure each model's proficiency in producing complete textual solutions unaided by external tools, as well as their competence in solving tasks through Python-based execution.

    DeepSeekMath-Base exhibits strong results on English benchmarks, rivaling the performance of the closed-source Minerva 540B \citep{minerva}, and outperforms every available open-source base model—such as Mistral 7B \citep{mistral} and Llemma-34B \citep{llemma}—irrespective of any prior math-specific training, often with a wide performance gap. Its advantages are even more pronounced on Chinese datasets, a result attributed to our divergence from preceding studies \citep{minerva,llemma}: unlike prior approaches that rely exclusively on English data for mathematical pre-training, we integrate diverse, high-quality datasets in other languages as well. After supplementing with mathematical instruction tuning and reinforcement learning, DeepSeekMath-Instruct and DeepSeekMath-RL set a new open-source standard, achieving over 50\% accuracy on the highly challenging, competition-grade MATH benchmark, a feat previously unseen in the open-source arena.
\end{itemize}

\item \textbf{Formal Mathematics}: DeepSeekMath-Base is assessed using the informal-to-formal theorem proving task introduced by \citep{dsp_proof} on the miniF2F benchmark \citep{minif2f}, employing Isabelle \citep{isabelle} as the proof assistant. The model displays robust few-shot autoformalization capabilities in this setting.
\item \textbf{Natural Language Understanding, Reasoning, and Code}: To thoroughly evaluate the general competencies of DeepSeekMath-Base in language comprehension, reasoning, and programming, we benchmark it on the Massive Multitask Language Understanding (MMLU) suite \citep{mmlu}, which features 57 multiple-choice tasks across a wide array of subjects. In addition, we use BIG-Bench Hard (BBH) \citep{bbh}, comprised of 23 tasks that predominantly demand intricate multi-step reasoning, alongside HumanEval \citep{codex} and MBPP \citep{mbpp}, two standard metrics for assessing code generation abilities. The results highlight that pre-training on mathematical data improves both general language understanding and reasoning proficiency.

\section{Math Pre-Training}

\subsection{Data Collection and Decontamination} This section describes the methodology for assembling the DeepSeekMath Corpus from Common Crawl. As illustrated in Figure \ref{fig:data_collect}, we developed an iterative pipeline designed to methodically extract a large-scale mathematics dataset from Common Crawl, beginning with a small but high-quality math-centric seed corpus. Notably, this framework extends readily to other domains, including coding.

To initiate the process, we utilize OpenWebMath \citep{openwebmath}, which offers a curated set of high-quality mathematical content from the web, as the foundational seed corpus. A fastText model \citep{joulin2016fasttext} is then trained to identify additional web pages that resemble the OpenWebMath dataset. For training, 500,000 examples are drawn at random from the seed as positive instances, while another 500,000 web pages from Common Crawl serve as negative counterparts. The model is trained using an open-source library\footnote{\url{https://fasttext.cc}}, with parameters set as follows: vector dimension at 256, learning rate of 0.1, a maximum word n-gram length of 3, minimum word occurrence threshold at 3, and 3 epochs of training. To compress the size of the initial Common Crawl dataset, we implement both URL deduplication and near-deduplication methods, thereby narrowing the data pool to 40 billion HTML pages. The fastText model is subsequently applied to these deduplicated pages to recall mathematical content. To ensure high data quality, the retrieved pages are ranked by their fastText-derived scores, and only those with the highest ranks are retained. The quantity of data maintained is empirically determined through pre-training runs on subsets containing the top 40B, 80B, 120B, and 160B tokens; for the initial iteration, we select the top 40B tokens for further use.

Following the initial round of data acquisition, a substantial portion of mathematical web pages remains uncollected, primarily due to the limited diversity within the positive examples used to train the fastText model. To address this, we incorporate additional mathematical web sources to expand and diversify the seed corpus, enhancing the performance of the fastText model. Concretely, we first segment the entire Common Crawl dataset into non-overlapping domains, where each domain represents a group of web pages sharing the same base URL. We then compute, for each domain, the proportion of its web pages included in the first collection pass. Domains where more than 10\% of their pages have been captured are flagged as math-related (for example, \url{mathoverflow.net}).

Next, we manually annotate URLs within these domains that contain mathematical content (such as \url{mathoverflow.net/questions}). Web pages connected to these annotated URLs that were missed in earlier collections are subsequently added to the seed corpus. This strategy augments our set of positive examples, facilitating the training of an improved fastText model that better retrieves mathematical data in future iterations. Ultimately, after completing four rounds of collection, we accumulate 35.5M mathematical web pages, representing a total of 120B tokens. In the fourth round, we observe that roughly 98\% of the dataset overlaps with the third iteration, prompting us to halt further data collection.

To mitigate the risk of contaminating benchmarks, we adopt the methodology of \cite{deepseek-coder}, removing any web pages containing questions or answers from established English mathematical benchmarks (e.g., GSM8K \citep{gsm8k}, MATH \citep{MATH}) as well as Chinese benchmarks (e.g., CMATH \citep{wei2023cmath}, AGIEval \citep{agieval}). The specific filtering process eliminates any text segment if it includes a 10-gram that is an exact match for any substring from the evaluation benchmarks. For benchmark passages that are between 3 and 9 grams in length, we apply exact matching to exclude affected web pages from the math training set.

\subsection{Assessing the DeepSeekMath~Corpus Quality}

To assess the effectiveness of the DeepSeekMath~Corpus, we conduct pre-training experiments comparing it against recently introduced math-specific corpora:

\begin{itemize}[topsep=0pt]
    \item \textbf{MathPile} \citep{mathpile}: an 8.9B token collection from various sources including textbooks, Wikipedia, ProofWiki, CommonCrawl, StackExchange, and arXiv, with arXiv making up over 85\% of the dataset;
    \item \textbf{OpenWebMath} \citep{openwebmath}: a corpus sourced from CommonCrawl, filtered for mathematical relevance, consisting of 13.6B tokens;
    \item \textbf{Proof-Pile-2} \citep{llemma}: comprising OpenWebMath, AlgebraicStack (containing 10.3B tokens of mathematical code), and 28.0B tokens from arXiv papers. For experiments using Proof-Pile-2, we adhere to the arXiv:Web:Code ratio of 2:4:1 as described in \cite{llemma}.
\end{itemize}

\subsubsection{Training Setting}
\label{sec:quality-policy}
We conduct mathematics-focused training using a 1.3B-parameter general-purpose language model, which adopts the architecture found in DeepSeek LLMs \citep{deepseek-llm}, hereafter referred to as DeepSeek-LLM 1.3B. For each mathematical dataset, we train the model independently over 150B tokens. All training procedures employ the lightweight and resource-efficient HAI-LLM framework \citep{haillm}. In line with the DeepSeek LLMs methodology, we utilize the AdamW optimizer \citep{adamW} configured with $\beta_1=0.9$, $\beta_2=0.95$, and a $\mathrm{weight\_decay}$ of $0.1$. The learning rate schedule follows a multi-phase approach: after a 2,000-step warmup to the maximum learning rate, the rate is linearly annealed to 31.6\% of its peak at 80\% of total training, then further reduced to 10.0\% by 90\% of completion. The peak learning rate is set at $5.3 \times 10^{-4}$, the batch size at 4M tokens, and context window at 4K tokens.

\subsubsection{Evaluation Results}

\textbf{The DeepSeekMath~Corpus exhibits exceptional quality, broad multilingual coverage, and stands as the most extensive resource in its class.}

\begin{itemize}[topsep=0pt]
    \item \textbf{High-quality}:
    To assess the downstream capabilities, we measure model accuracy across eight math benchmarks using few-shot chain-of-thought prompting \cite{cot}. Table \ref{tab:corpora_comparison} reveals a substantial performance advantage for models trained on DeepSeekMath~Corpus. Additionally, Figure \ref{fig:corpora_comparisons} illustrates that DeepSeekMath-trained models surpass those trained on Proof-Pile-2 after 50B tokens (equivalent to a single epoch over Proof-Pile-2), suggesting that DeepSeekMath~Corpus delivers superior average data quality.
    \item \textbf{Multilingual}:
    DeepSeekMath~Corpus contains mathematics texts spanning several languages, with English and Chinese being most prominent. Evidence from Table \ref{tab:corpora_comparison} demonstrates that this corpus improves mathematical reasoning abilities in both English and Chinese. By contrast, traditional mathematical corpora—largely restricted to English—yield only marginal benefits for Chinese, sometimes even diminishing its performance in that domain.
    \item \textbf{Large-scale}:
    Compared to existing mathematical datasets, the DeepSeekMath~Corpus is significantly larger. As presented in Figure \ref{fig:corpora_comparisons}, training DeepSeek-LLM 1.3B on this corpus leads to a sharper and more sustained rise in performance, whereas models using smaller, baseline corpora—already cycled through several epochs—tend to quickly stagnate in terms of accuracy improvements.
\end{itemize}

\subsection{Training and Evaluation of DeepSeekMath-Base 7B}

This section details the development of DeepSeekMath-Base 7B, a foundational model characterized by robust mathematical reasoning capabilities. The initialization is based on DeepSeek-Coder-Base-v1.5 7B \citep{deepseek-coder}, and training proceeds for a total of 500B tokens. The data utilized for training are distributed as follows: 56\% originates from the DeepSeekMath Corpus, 4\% from AlgebraicStack, 10\% from arXiv, 20\% from Github code repositories, and the remaining 10\% consists of natural language content drawn from Common Crawl in both English and Chinese. The primary training configuration adheres to the guidelines provided in Section \ref{sec:quality-policy}, with two modifications: the peak learning rate is increased to 4.2e-4, and the batch size is set to 10 million tokens.

A thorough evaluation is carried out to analyze the mathematical proficiency of DeepSeekMath-Base 7B. The analysis encompasses its effectiveness in generating self-contained mathematical solutions without tool assistance, its performance when leveraging external tools for mathematical tasks, and its competency in formal theorem verification. In addition, we assess general abilities beyond mathematics, specifically examining natural language understanding, reasoning, and programming proficiency.

\paragraph{Stepwise Mathematical Problem Solving}

To measure DeepSeekMath-Base's problem-solving abilities in mathematics, we employ few-shot chain-of-thought prompting \citep{cot} on eight distinct benchmarks in both English and Chinese. These benchmarks address quantitative reasoning tasks (for example, GSM8K \citep{gsm8k}, MATH \citep{MATH}, and CMATH \citep{wei2023cmath}) as well as multiple-choice assessments (such as MMLU-STEM \citep{mmlu} and Gaokao-MathQA \citep{agieval}), thereby encompassing a wide spectrum of mathematical disciplines ranging from elementary to undergraduate level difficulty.

As detailed in Table \ref{tab:base_math_cot}, DeepSeekMath-Base 7B consistently achieves top results on all eight benchmarks among open-source base models—surpassing widely adopted models such as Mistral 7B \citep{mistral} and the more recent Llemma 34B \citep{llemma}, the latter of which received targeted mathematical pretraining on Proof-Pile-2 \citep{llemma}. Particularly remarkable is its performance on the MATH competition-level dataset, where DeepSeekMath-Base exceeds the scores of other open-source base models by more than 10 percentage points in absolute terms, and also outperforms Minerva 540B \citep{minerva}—a proprietary model based on PaLM \citep{palm} and 77 times larger—which was further refined on mathematical literature.

\paragraph{Mathematical Problem Solving with Tool Use} We assess the models’ ability to conduct program-assisted mathematical reasoning on GSM8K and MATH via few-shot program-of-thought prompting \citep{pot,pal}. In this setting, the model is tasked with solving each problem by generating a Python script, permitted to use libraries such as \textit{math} and \textit{sympy} for handling complex calculations. The answer is taken as the output produced by running the script. According to Table \ref{tab:base_math_pot_proof}, DeepSeekMath-Base 7B surpasses the former state-of-the-art, Llemma 34B, in these tasks.

\paragraph{Formal Mathematics}

Automating formal proofs is instrumental in ensuring correctness and improving productivity, a topic gaining growing interest. We test DeepSeekMath-Base 7B’s proficiency on informal-to-formal proof tasks as described in \citep{dsp_proof}, where the goal is to create a formalized proof based on an informal problem statement, its formal expression, and an informal argument. The evaluation uses miniF2F \citep{minif2f}, which benchmarks formal solutions for Olympiad-level mathematics. For each problem, formal proofs in Isabelle are generated using few-shot prompting. Consistent with \cite{dsp_proof}, models are tasked with producing proof sketches, after which the automated prover Sledgehammer \citep{sledgehammer} completes any missing proof steps. Table \ref{tab:base_math_pot_proof} demonstrates that DeepSeekMath-Base 7B exhibits robust capability in automating proof formalization.

\paragraph{Natural Language Understanding, Reasoning, and Code}

To measure proficiency in language understanding, reasoning, and programming, we evaluate the models on MMLU \citep{mmlu} (for comprehension), BBH \citep{bbh} (for reasoning), and on the HumanEval \citep{codex} and MBPP \citep{mbpp} code generation benchmarks. Table \ref{tab:understanding_reasoning_code} shows that DeepSeekMath-Base 7B substantially improves upon the performance of its predecessor, DeepSeek-Coder-Base-v1.5 \citep{deepseek-coder}, on MMLU and BBH, highlighting the effectiveness of mathematical training in bolstering language and reasoning abilities. Moreover, by incorporating coding tokens during continual training, DeepSeekMath-Base 7B successfully retains the coding benchmark performance of DeepSeek-Coder-Base-v1.5. Across these three reasoning and coding tasks, DeepSeekMath-Base 7B achieves significantly better results than the generalist Mistral 7B \citep{mistral}.

\section{Supervised Fine-Tuning}

\subsection{SFT Data Curation}

We assemble a dataset designed for mathematical instruction tuning that encompasses both English and Chinese questions, spanning a broad spectrum of mathematical topics and complexity levels. Each problem is matched with a solution articulated in chain-of-thought (CoT) \citep{cot}, program-of-thought (PoT) \citep{pot,pal}, or tool-integrated reasoning style \citep{tora}. In total, the training corpus comprises 776K samples.

\begin{itemize}[topsep=0pt]
    \item \textbf{English mathematical datasets}: We enrich GSM8K and MATH tasks with annotations that employ tool-integrated reasoning. Additionally, we include a curated subset of MathInstruct \citep{MathInstruct} and leverage the training set from Lila-OOD \citep{lila}, which features problems resolved using either CoT or PoT strategies. This English collection spans a wide range of mathematical domains such as algebra, probability, number theory, calculus, and geometry.
    \item \textbf{Chinese mathematical datasets}: Our collection of Chinese problems targets K-12 mathematics, encompassing 76 distinct sub-topics, including, for example, linear equations. Solutions are systematically provided in both CoT and tool-integrated reasoning formats.
\end{itemize}

\subsection{Training and Evaluating DeepSeekMath-Instruct 7B}

Here, we present the DeepSeekMath-Instruct 7B model, which is fine-tuned on mathematical instructions using the DeepSeekMath-Base as the foundation. During training, input samples are concatenated at random, up to a maximum sequence length of 4K tokens. The fine-tuning process is conducted over 500 training steps, utilizing a batch size of 256, with a fixed learning rate set at 5e-5.

To assess mathematical proficiency, we test models in both tool-augmented and tool-free scenarios, using four benchmarks for quantitative reasoning tasks in English and Chinese. For context, we compare our approach with the strongest contemporaneous models:

\begin{itemize}[topsep=0pt]
    \item \textbf{Closed-source models} are comprised of: (1) the GPT series, including GPT-4 \citep{gpt4} and the advanced GPT-4 Code Interpreter~\footnote{\url{https://openai.com/blog/chatgpt-plugins\#\#code-interpreter}}; (2) Gemini Ultra and Pro \citep{gemini}; (3) Inflection-2 \citep{inflection-2}; (4) Grok-1~\footnote{\url{https://x.ai/model-card}}, as well as newly released Chinese models such as (5) Baichuan-3~\footnote{\url{https://www.baichuan-ai.com}}, and (6) the most recent GLM-4~\footnote{\url{https://open.bigmodel.cn/dev/api\#glm-4}}

    from the GLM model family \citep{glm}.
\end{itemize}

These models are designed for broad applicability, and most have been subjected to various alignment strategies.

\item \textbf{Open-source models} considered in this study include: foundational models such as (1) DeepSeek-LLM-Chat 67B \citep{deepseek-llm}, (2) Qwen 72B \citep{qwen}, (3) SeaLLM-v2 7B \citep{seallm}, and (4) ChatGLM3 6B \citep{chatglm3}. Additionally, several models incorporate specialized mathematical capabilities, including (5) InternLM2-Math 20B~\footnote{\url{https://github.com/InternLM/InternLM-Math}}, which extends InternLM2 via math-centric pretraining followed by instruction-based finetuning, (6) Math-Shepherd-Mistral 7B, which employs PPO training \citep{schulman2017proximal} on top of Mistral 7B \citep{mistral} using a process-supervised reward mechanism, (7) the WizardMath series \citep{wizardmath}, where mathematical reasoning is bolstered in Mistral 7B and Llama-2 70B \citep{llama2} using evolve-instruct—a variant of instruction finetuning based on AI-generated instructions—and PPO, with training problems largely drawn from GSM8K and MATH, (8) MetaMath 70B \citep{metamath}, constructed by further finetuning Llama-2 70B on an augmented GSM8K and MATH dataset, (9) ToRA 34B \cite{tora}, which adapts CodeLlama 34B for tool-augmented mathematical reasoning, and (10) MAmmoTH 70B \citep{MathInstruct}, achieved through instruction finetuning of Llama-2 70B on MathInstruct.

Table \ref{tab:sft_rl_math} displays results in the setting where external tools are prohibited: here, DeepSeekMath-Instruct 7B achieves remarkable stepwise reasoning capabilities. In particular, on the challenging MATH benchmark, our model outperforms all other open-source systems and most commercial counterparts (including Inflection-2 and Gemini Pro) by a margin of at least 9 percentage points. This advantage is observed even against considerably larger models (such as Qwen 72B) or those specifically optimized for mathematical reasoning via reinforcement learning (such as WizardMath-v1.1 7B). While DeepSeekMath-Instruct's results approach those of leading Chinese closed-source models like GLM-4 and Baichuan-3 on MATH, it remains behind GPT-4 and Gemini Ultra.

When models are permitted to incorporate both natural language reasoning and code-based tools in their solutions, DeepSeekMath-Instruct 7B reaches nearly 60\% accuracy on MATH, establishing itself as the best among open-source models under this protocol. Across other benchmarks, its performance is on par with DeepSeek-LLM-Chat 67B—the former top performer among open-source systems, despite being an order of magnitude larger.

\section{Reinforcement Learning}

\subsection{Group Relative Policy Optimization}
Reinforcement learning (RL) has been shown to significantly boost LLM mathematical reasoning following Supervised Fine-Tuning (SFT) \citep{wang2023math,wizardmath}. Here, we propose Group Relative Policy Optimization (GRPO), our resource-efficient RL framework.

\subsubsection{From PPO to GRPO} Proximal Policy Optimization (PPO) \citep{schulman2017proximal} is an actor-critic RL technique that has become standard for RL-based finetuning of LLMs \citep{ouyang2022training}. PPO improves model policies by maximizing the following surrogate loss:

\begin{equation}
\footnotesize
    \mathcal{J}_{PPO}(\theta) = \mathbb{E}{[q \sim P(Q), o \sim \pi_{\theta_{old}}(O|q)]} \frac{1}{|o|} \sum_{t=1}^{|o|} \min \left[ \frac{\pi_\theta(o_{t} | q, o_{<t})}{\pi_{\theta_{old}}(o_{t} | q, o_{<t})} A_{t}, \text{clip} \left( \frac{\pi_\theta(o_{t} | q, o_{<t})}{\pi_{\theta_{old}}(o_{t} | q, o_{<t})}, 1 - \varepsilon, 1 + \varepsilon \right)  A

Here, $\pi_{\theta}$ denotes the current policy, while $\pi_{\theta_{old}}$ refers to the previous iteration of the policy model. The variables $q$ and $o$ represent questions and sampled outputs, respectively, drawn from a dataset of questions and from the distribution defined by $\pi_{\theta_{old}}$. The parameter $\varepsilon$ serves as the clipping threshold, an essential hyper-parameter in PPO that stabilizes policy updates. The term $A_t$ stands for the advantage estimate, computed via Generalized Advantage Estimation (GAE) \citep{gae} utilizing the future rewards $\{r_{\ge t}\}$ in conjunction with a learned value function $V_{\psi}$. Consequently, the PPO algorithm necessitates concurrent training of a value function alongside the policy. To avoid excessive optimization towards the reward model, it is common to introduce a per-token KL penalty, calculated against a reference model, within the reward signal for each token, as prescribed in \citep{ouyang2022training}:

\begin{equation}
    r_{t} = r_\varphi(q, o_{\le t}) - \beta \log\frac{\pi_{\theta}(o_{t}|q, o_{<t})}{\pi_{ref}(o_{t}|q, o_{<t})},
\label{eq:PPO-reward}
\end{equation}

Here, $r_\varphi$ represents the reward function, $\pi_{ref}$ is a reference policy (commonly, the initial SFT model), and $\beta$ regulates the magnitude of the KL penalty.

The use of a value function in PPO often necessitates a neural network with complexity comparable to that of the policy, thereby incurring significant memory and computational overhead. Additionally, within RL fine-tuning, the value function serves primarily as a baseline to lower the variance of the advantage estimator. In the context of LLM training, the reward model typically produces nonzero rewards only for the last token, which complicates the reliable supervision of the value function for every token. To remedy these issues, we introduce Group Relative Policy Optimization (GRPO), illustrated in Figure \ref{fig:grpo}, which eliminates the need for explicit value function approximation as in PPO. GRPO instead employs the mean reward across a collection of output samples—each conditioned on the same question—as a baseline. In practice, for a given question $q$, a batch of outputs $\{o_1, o_2, \cdots, o_G\}$ is drawn from $\pi_{\theta_{old}}$. The policy is then updated by maximizing the objective:

\begin{equation}
\footnotesize
\begin{split}
    \mathcal{J}_{GRPO}(\theta) &= \mathbb{E}{[q \sim P(Q), \{o_i\}_{i=1}^G \sim \pi_{\theta_{old}}(O|q)]}  \\
    & \frac{1}{G}\sum_{i=1}^G\frac{1}{|o_i|} \sum_{t=1}^{|o_i|} \left\{ \min \left[ \frac{\pi_\theta(o_{i,t} | q, o_{i,<t})}{\pi_{\theta_{old}}(o_{i,t} | q, o_{i,<t})} \hat{A}_{i,t}, \text{clip} \left( \frac{\pi_\theta(o_{i,t} | q, o_{i,<t})}{\pi_{\theta_{old}}(o_{i,t} | q, o_{i,<t})}, 1 - \varepsilon, 1 + \varepsilon \right)  \hat{A}_{i,t} \right] - \beta \mathbb{D}_{KL}\left[\pi_{\theta} || \pi_{ref}\right]\right\} ,
\end{split}
\label{eq:GRPO-obj}
\end{equation}

Here, $\varepsilon$ and $\beta$ remain as hyper-parameters. The term $\hat{A}_{i,t}$ is an advantage estimate computed by contrasting rewards only within the group of sampled outputs for a particular question. Details regarding this group-relative computation will be elaborated in subsequent sections. By designing the advantage calculation to be relative within a group, GRPO aligns with the comparative structure of reward models, which are typically trained on datasets containing pairwise or setwise preference information for the same prompt. Additionally, unlike PPO where the KL term is included in the reward, GRPO incorporates a KL divergence regularizer directly into the loss function, thereby simplifying the computation of $\hat{A}_{i,t}$. Distinct from the KL penalty applied in (\

\begin{algorithm}[t]
  \small
  \caption{Iterative Group Relative Policy Optimization}
  \textbf{Input} initial policy model $\pi_{\theta_{\text{init}}}$; reward models $r_\varphi$; task prompts $\mathcal{D}$; 
  hyperparameters $\varepsilon$, $\beta$, $\mu$
  \begin{algorithmic}[1]
    \State Initialize policy model: $\pi_\theta \leftarrow \pi_{\theta_{\text{init}}}$
    \For{iteration = 1, \dots, I}
       \State Set the reference model $\pi_{ref} \leftarrow \pi_{\theta}$
      \For{step = 1, \dots, M}
      \State Select a minibatch $\mathcal{D}_b$ from $\mathcal{D}$
      \State Update previous policy: $\pi_{\theta_{old}} \leftarrow \pi_{\theta}$ 
      \State For each question $q \in \mathcal{D}_b$, sample $G$ responses $\{o_i\}_{i=1}^G$ using $\pi_{\theta_{old}} (\cdot \mid q) $
      \State Use $r_\varphi$ to evaluate each $o_i$ and obtain rewards $\{r_i\}_{i=1}^{G}$
      \State For every output $o_i$, compute the advantage $\hat{A}_{i,t}$ for the $t$-th token through group relative advantage calculation
      \For{GRPO iteration = 1, \dots, $\mu$}
        \State Adjust policy $\pi_{\theta}$ by optimizing the GRPO criterion (Equation \ref{eq:GC-GRPO})
      \EndFor
    \EndFor
    \State Periodically refine $r_\varphi$ via ongoing training with a replay buffer
    \EndFor 
  \end{algorithmic}
  \textbf{Output} $\pi_\theta$
  \label{alg:iter-grpo}
\end{algorithm}

\subsubsection{Outcome Supervision RL with GRPO} To elaborate, for any given question $q$, a collection of $G$ candidate outputs $\{o_1, o_2, \cdots, o_G\}$ are generated by sampling from the old policy $\pi_{\theta_{old}}$. Each candidate response is evaluated by a reward model, yielding a set of rewards $\mathbf{r} = \{r_1, r_2, \cdots, r_G\}$. The set of rewards is normalized by removing the mean and scaling by the standard deviation within the group. In the context of outcome supervision, this normalized terminal reward is assigned as the advantage $\hat{A}_{i, t}$ at every token position $t$ of $o_i$, specifically $\hat{A}_{i, t} = \widetilde{r}_i = \frac{r_i- {\rm mean}(\mathbf{r})}{{\rm std}(\mathbf{r})}$. The policy is then improved by maximizing the objective from equation (\ref{eq:GRPO-obj}).

\subsubsection{Process Supervision RL with GRPO} While outcome supervision restricts feedback to a single reward for each complete output, which can be limiting in supervising policies on intricate mathematical tasks, process supervision—as in \cite{wang2023math}—offers a more granular reward signal by scoring each intermediate reasoning step. In this formulation, given a question $q$ and sampled outputs $\{o_1, o_2, \cdots, o_G\}$, a process-level reward model produces a set of stepwise rewards for each output, denoted as $\mathbf{R} = \{ \{r_1^{index(1)},\cdots,r_1^{index(K_1)}\}, \cdots,  \{r_G^{index(1)},\cdots,r_G^{index(K_G)}\} \}$, with $index(j)$ indicating the token position where the $j$-th step ends, and $K_i$ being the number of reasoning steps in $o_i$. These step rewards are also normalized: $\widetilde{r}_i^{index(j)} = \frac{r_i^{index(j)} - {\rm mean(\mathbf{R})}}{{\rm std(\mathbf{R})}}$. The advantage for each token is computed as the cumulative normalized reward over the steps from the current token onward, expressed as $\hat{A}_{i, t} = \sum_{index(j) \ge t} \widetilde{r}_i^{index(j)}$. The policy is subsequently refined by maximizing the same GRPO

\subsubsection{Iterative RL with GRPO} As reinforcement learning advances, previously trained reward models may become inadequate for effectively guiding the continually updated policy model. In response, we implement an iterative RL approach utilizing GRPO. Detailed in Algorithm \ref{alg:iter-grpo}, this methodology entails generating new datasets for the reward model through samples produced by the evolving policy model. We then update the reward model using a replay buffer that preserves 10\% of data from prior iterations. The latest policy model is adopted as the reference for subsequent training phases, during which both the reward and policy models are jointly refined, leveraging newly obtained reward data.

\subsection{Training and Evaluation of DeepSeekMath-RL}

Our RL experiments are carried out on the DeepSeekMath-Instruct 7B model. The RL training dataset comprises approximately 144K chain-of-thought-style queries derived from the GSM8K and MATH sections within the SFT data. We deliberately exclude other SFT-derived questions to assess RL performance specifically on evaluation sets devoid of their examples during RL optimization. The reward model training data is curated in accordance with procedures outlined by \citep{wang2023math}. Initialization of the reward model is performed on DeepSeekMath-Base 7B with a learning rate set at 2e-5. The policy model, under GRPO, is optimized with a learning rate of 1e-6, and the KL divergence coefficient is fixed at 0.04. For each prompt, 64 samples are generated, with a maximum token length capped at 1024 and a batch size of 1024 for training. Each iteration consists of a single policy update following exploration. DeepSeekMath-RL 7B’s evaluation follows the same benchmarks as DeepSeekMath-Instruct 7B. Within DeepSeekMath-RL 7B, tasks based on GSM8K and MATH that require chain-of-thought reasoning are classified as in-domain, while all other benchmark tasks are treated as out-of-domain.

Table \ref{tab:sft_rl_math} presents comparative results for both open-source and proprietary models, evaluating their performance using chain-of-thought and tool-augmented reasoning approaches across English and Chinese datasets. The main observations are as follows: (1) When leveraging chain-of-thought prompting, DeepSeekMath-RL 7B achieves 88.2\% accuracy on GSM8K and 51.7\% on MATH, exceeding the scores obtained by all open-source models ranging from 7B to 70B parameters, as well as surpassing most closed-source alternatives. (2) Significantly, the only additional training data provided to DeepSeekMath-RL 7B are chain-of-thought-based instruction tuning examples from GSM8K and MATH, continuing from the DeepSeekMath-Instruct 7B checkpoint. Even with this limited instruction tuning scope, DeepSeekMath-RL 7B achieves higher results across every metric compared to DeepSeekMath-Instruct 7B, underscoring the benefit of reinforcement learning.

\section{Discussion}
This section summarizes insights derived from our experiments in both pre-training and reinforcement learning.

\subsection{Lessons Learnt in Pre-Training}

Here, we detail insights gathered from the pre-training phase. Unless noted otherwise, all experiments follow the setup described in Section \ref{sec:quality-policy}. In this context, references to the DeepSeekMath Corpus denote the 89B-token collection amassed during the second data iteration.

\subsubsection{Code Training Benefits Mathematical Reasoning}

A widely held but inadequately tested assumption posits that code pre-training enhances reasoning capabilities. Our findings offer evidence supporting this, specifically within mathematical tasks: code-based training consistently boosts models’ proficiency in mathematical reasoning, irrespective of whether external tools are used.

To probe the effects of code pre-training on mathematical reasoning, we designed both a two-stage and a single-stage training protocol:

\textbf{Two-Stage Training}

\begin{itemize}[topsep=0pt]
    \item \textbf{Code Training for 400B Tokens $\rightarrow$ Math Training for 150B Tokens}: We first expose DeepSeek-LLM 1.3B to 400 billion code tokens, followed by further training on 150 billion math-specific tokens.
    \item \textbf{General Training for 400B Tokens $\rightarrow$ Math Training for 150B Tokens}: For comparative purposes, we also conduct an experiment where, in the initial phase, general language data (sourced from a comprehensive general dataset curated by DeepSeek-AI) is used instead of code tokens. This is to determine the relative contributions of code versus general-domain tokens for improving mathematical problem-solving skills.
\end{itemize}

\textbf{One-Stage Training}

\begin{itemize}[topsep=0pt]
    \item \textbf{Math Training for 150B Tokens}: DeepSeek-LLM 1.3B is directly trained on 150 billion math-related tokens without prior exposure to code or general tokens.
    \item \textbf{Training on a mixture of 400B Code Tokens and 150B Math Tokens}: Because a sequential approach—math training subsequent to code exposure—tends to impair code-related capabilities, we evaluate whether a simultaneous, blended curriculum comprising both 400 billion code tokens and 150 billion math tokens enhances mathematical reasoning and prevents the detrimental effects of catastrophic forgetting.
\end{itemize}

\paragraph{Results}
The outcomes, as summarized in Table \ref{tab:code-math} and Table \ref{tab:code-math-general-eval}, reveal performance differences across various training paradigms.

Incorporating code data significantly advances program-assisted mathematical reasoning under both staged and combined training regimes. In the two-stage format, code-focused training alone notably improves DeepSeek-LLM 1.3B's proficiency at tackling GSM8K and MATH tasks with Python. Further gains are realized with the addition of math-centric training. Moreover, when code and math data are merged in a single-stage training procedure, the risks of catastrophic forgetting typical in two-phase pipelines are alleviated, while code generation skills (see Table \ref{tab:code-math-general-eval}) and program-aided mathematical problem-solving (see Table \ref{tab:code-math}) both benefit.

Exposure to code data also bolsters mathematical reasoning in settings without explicit tool usage. Sequentially introducing code tokens leads to moderate yet meaningful performance gains and enhances the impact of subsequent math-targeted instruction, ultimately producing the highest accuracy. Conversely, one-stage training with both code and math data appears to undermine math problem-solving in the absence of tool support. This could be attributed to the limited model size of DeepSeek-LLM 1.3B, which may constrain its ability to effectively encode and utilize both code and math information in tandem.

\subsubsection{ArXiv Papers Appear Unhelpful for Enhancing Mathematical Reasoning}

Despite their widespread inclusion as a core element in math pre-training datasets \citep{minerva,gpt-f,llemma,mathpile}, the specific influence of arXiv papers on mathematical reasoning remains under-examined. Somewhat unexpectedly, our results suggest that arXiv papers do not confer observable gains in mathematical reasoning ability. We assess this by training language models of varied capacities—including DeepSeek-LLM 1.3B and DeepSeek-Coder-Base-v1.5 7B \citep{deepseek-coder}—on arXiv data sets processed through distinct methods:

\begin{itemize}[topsep=0pt]
    \item \textbf{MathPile} \citep{mathpile}: An 8.9-billion-token resource constructed through extensive filtering and cleaning, where over 85\% of the content originates from scientific arXiv articles;
    \item \textbf{ArXiv-RedPajama} \citep{redpajama}: A 28.0-billion-token compilation of raw arXiv LaTeX documents, with ancillary components such as preambles, macros, comments, and bibliographies excised.
\end{itemize}

For our evaluation, we conduct independent training runs: DeepSeek-LLM 1.3B on 150B tokens and DeepSeek-Coder-Base-v1.5 7B on 40B tokens for each arXiv-based dataset. Results indicate a lack of effectiveness—when models are trained solely on arXiv data, there are neither measurable advancements nor sometimes even regressions in performance across the studied mathematical reasoning benchmarks. These encompass a range of difficulties and formats, such as GSM8K and MATH for quantitative reasoning (refer to Table \ref{tab:arxiv-cot}), multiple-choice assessments like MMLU-STEM (see Table \ref{tab:arxiv-cot}), and formal mathematics tests such as miniF2F (refer to Table \ref{tab:arxiv_atp}).

Nonetheless, these findings have important caveats and should not be regarded as definitive. Notably, we have not yet analyzed:

\begin{itemize}[topsep=0pt]
    \item The influence of arXiv data on math tasks not represented in our experiments, for instance, converting formal mathematical proofs to informal descriptions;
    \item The potential effects that arXiv data might exert when it is blended with additional training corpora;
    \item Whether performance improvements would emerge with further scaling of model parameters.
\end{itemize}

Consequently, these open questions call for more comprehensive study, which we defer to future investigations.

\subsection{Perspectives on Reinforcement Learning}
\subsubsection{In Pursuit of a Unified Framework} This section proposes a coherent analytical framework encompassing various training strategies, including SFT, RFT, DPO, PPO, and GRPO, and presents experiments to scrutinize factors within this framework. Typically, the gradient of the model parameter $\theta$ for a given method can be formulated as:

\begin{equation}
    \nabla_{\theta}\mathcal{J}_{\textcolor{red}{\mathcal{A}}}(\theta) = \mathbb{E}[\underbrace{(q,o) \sim \textcolor{red}{\mathcal{D}}}_{Data \ Source}]\left( \frac{1}{|o|} \sum_{t=1}^{|o|}  \underbrace{GC_{{\mathcal{A}}}(q, o, t, \textcolor{red}{\pi_{{rf}}})}_{Gradient \ Coefficient}  \nabla_{\theta}\log \pi_{\theta}(o_t | q, o_{<t})\right).
\label{eq:objective}
\end{equation}

This formulation highlights three fundamental elements: 1) \textit{Data Source $\mathcal{D}$}, defining the dataset for model training; 2) \textit{Reward Function $\pi_{{rf}}$}, responsible for generating the reward signals guiding training; and 3) \textit{Algorithm $\mathcal{A}$}, which utilizes the training samples and reward feedback to calculate the gradient coefficient $GC$, thereby modulating the learning signal as either encouragement or penalization. We examine multiple commonly adopted approaches within this shared framework:

\begin{itemize}
    \item  \textbf{Supervised Fine-tuning (SFT)}: In SFT, the pretrained model is adapted using curated SFT datasets generated by human annotators.
    \item \textbf{Rejection Sampling Fine-tuning (RFT)}: RFT continues to train the SFT model using outputs sampled from the SFT model, filtered according to the correctness of the generated answers to SFT prompts.
    \item \textbf{Direct Preference Optimization (DPO)}: DPO builds on the SFT model, tuning it further using additional outputs drawn from the SFT model and employing a pairwise DPO objective.
    \item \textbf{Online Rejection Sampling Fine-tuning (Online RFT)}: Unlike standard RFT, Online RFT initializes the policy network from SFT, but iteratively updates it with new outputs sampled from the most recent version of the policy itself.
    \item \textbf{PPO/GRPO}: In PPO/GRPO, the model is initialized with SFT weights and subsequently optimized via reinforcement on outputs sampled from the ongoing policy.
\end{itemize}

A breakdown of the constituent elements of these methods is provided in Table \ref{tab:data_policy}. For an in-depth exploration of the derivation procedures, please consult Appendix \ref{app:analysis-rl}.

\paragraph{Reflection on Data Source} The origin of the data is categorized into two types: online sampling and offline sampling. Online sampling refers to collecting training data generated from the policy model’s behavior during real-time exploration, whereas offline sampling corresponds to data obtained from the outputs of the pretrained SFT model. In this context, RFT and DPO utilize offline data, while Online RFT and GRPO rely on online data. 

Referencing Figure \ref{fig:rl-analysis},
it is evident that Online RFT delivers notably superior performance compared to RFT on both benchmarks examined. During the early phases of training, Online RFT and RFT yield similar results; however, as training progresses, Online RFT establishes a marked lead, illustrating the clear benefits of online data acquisition. This finding is consistent with expectations: early in training, the actor and the SFT model behave similarly, so their sampled outputs show minimal divergence. Later in training, the actor’s sampled data increasingly diverges from the initial SFT outputs, making dynamic, real-time sampling more effective.

\paragraph{Reflection on Gradient Coefficient} The procedure by which the reward is translated into a gradient coefficient for updating model parameters varies. In our experiments, the reward function is operationalized as either a ‘Rule’—where response quality is assessed by factual accuracy—or as a ‘Model’—where a reward model assigns scores to responses, trained on rule-based assessments. As detailed in Equations \ref{eq:GC-RFT} and \ref{eq:GC-GRPO}, GRPO and Online RFT diverge in how they incorporate the reward model: GRPO’s gradient coefficient is uniquely modulated in proportion to the reward value assigned, enabling differentiated strengthening or weakening of responses based on their relative merit. Online RFT, by contrast, does not provide additional penalization for erroneous responses and applies the same degree of reinforcement to all correctly answered samples.

As illustrated in Figure \ref{fig:rl-analysis}, GRPO outperforms online RFT, underscoring the effectiveness of modifying positive and negative gradient coefficients. Moreover, GRPO+PS achieves better results than GRPO+OS, demonstrating the advantage of implementing step-aware, fine-grained gradient coefficients. We also examine iterative RL by carrying out two rounds of optimization in our experiments. Results in Figure \ref{fig:iter-rl} reveal that this iterative approach notably boosts performance, particularly during the initial iteration.

\subsubsection{Why RL Works?} This work applies reinforcement learning to a selected portion of instruction tuning data, achieving marked improvements relative to the baseline instruction tuning model. To further elucidate the mechanisms behind RL’s efficacy, we compare Pass@K and Maj@K accuracies of the Instruct and RL-enhanced models across two benchmarks. According to Figure \ref{fig:maj-pass}, RL mainly increases Maj@K accuracy while Pass@K remains largely unchanged. This observation suggests that RL strengthens overall model performance primarily by enhancing the resilience of the output distribution; \textbf{in essence, the gain appears to come from elevating the likelihood of selecting the correct answer from the TopK candidates, rather than improving the intrinsic ability of the model.} In a similar vein, \citep{wang2023making} noted a \textbf{misalignment problem} affecting reasoning tasks in SFT models, and documented that various preference alignment strategies can remediate this issue and bolster reasoning accuracy \citep{yuan2023rrhf,song2023preference,wang2023making}.

\subsubsection{How to Achieve More Effective RL?}
\label{sec:effec_rl}
Our results verify that RL is highly effective in mathematical reasoning contexts. Additionally, we introduce a unified framework that encapsulates various prominent training strategies under the broader view of direct or reduced RL approaches. This framework, detailed in Equation \ref{eq:objective}, consists of three principal elements: Data Source, Algorithm, and Reward Function. We discuss several promising directions for advancing each of these aspects.

\paragraph{Data Source} The data source forms the foundational input for any training strategy. For RL, we define the data source as the pool of unlabeled questions along with candidate outputs generated from the policy model. In our study, we exclusively utilize instruction tuning questions and employ a straightforward nucleus sampling method for output generation. We believe that this design choice could partly explain why our RL framework primarily enhances Maj@K. Looking ahead, we plan to evaluate our RL process with out-of-distribution prompts and integrate \textbf{advanced sampling (decoding) strategies}, such as those leveraging tree-search algorithms \citep{yao2023tree}. Moreover, incorporating \textbf{efficient inference techniques} \citep{xia-etal-2023-speculative,leviathan2023fast,kwon2023efficient,xia2024unlocking}, which are crucial for the exploration effectiveness of policy models, is another important avenue for further improvement.

\paragraph{Algorithms} Algorithms utilize both data and reward feedback to adjust the gradient coefficients and thus update model parameters. As delineated in Equation \ref{eq:objective}, contemporary methods generally \textbf{TRUST} the reward function’s signal, using it to modulate the conditional probability of specific tokens either upwards or downwards. Nonetheless, the reliability of the reward signal cannot be guaranteed in all scenarios, particularly in tasks characterized by high complexity. To illustrate, even datasets like PRM800K \citep{lightman2023let}, which have undergone meticulous annotation by highly trained personnel, exhibit an error rate close to 20\%\footnote{\url{https://github.com/openai/prm800k/issues/12\#issuecomment-1728491852}} in labeling. This observation motivates our investigation into reinforcement learning algorithms with improved robustness to noisy reward signals. We contend that \textbf{WEAK-TO-STRONG} alignment approaches \citep{burns2023weak} represent a potentially transformative direction for algorithmic development.

\paragraph{Reward Function} The reward function underpins the entire training process by delivering the feedback signal. In the context of RL, this function is typically implemented as a neural reward model. We propose three key avenues for advancing reward models: 1) \textbf{Strengthening the reward model’s generalization capacity.} The reward model should be capable of robust performance when confronted with out-of-distribution queries and sophisticated decoded responses; failure here means that reinforcement learning will only maintain the status quo of LLMs’ output distribution without substantive progress in their abilities. 2) \textbf{Capturing and utilizing the reward model’s uncertainty.} Modeling uncertainty could be instrumental in connecting less reliable reward models with weak-to-strong algorithmic frameworks. 3) \textbf{Constructing effective, high-fidelity process reward models} capable of furnishing detailed and nuanced feedback throughout the reasoning process \citep{lightman2023let,wang2023math}.

\section{Conclusion, Limitation, and Future Work}

This work introduces DeepSeekMath, a model that sets a new benchmark among open-source systems for competition-level MATH tasks and narrows the performance gap with proprietary counterparts. DeepSeekMath~builds on DeepSeek-Coder-v1.5 7B and is subjected to sustained training over 500B tokens, prominently featuring 120B mathematical tokens curated from Common Crawl within its dataset. Our thorough ablation investigations reveal that web-based data serves as a rich resource for quality mathematical content, whereas arXiv contributions did not yield the anticipated benefits. We propose Group Relative Policy Optimization (GRPO), a modified form of Proximal Policy Optimization (PPO), which enhances mathematical reasoning abilities while demanding less memory overhead. Empirical findings indicate that GRPO maintains efficacy, even when DeepSeekMath-Instruct 7B attains elevated benchmark scores. Additionally, we outline a unified framework to conceptualize related approaches and highlight several promising avenues for advancing reinforcement learning techniques.

Despite DeepSeekMath's~notable achievements on quantitative reasoning assessments, it still lags behind closed-source models in domains such as geometry and formal theorem-proving. As observed in preliminary evaluations, the model fails to resolve tasks involving triangles and ellipses, which likely points to underlying biases in the selection of pre-training and fine-tuning data. Moreover, due to its limited scale, DeepSeekMath~is surpassed by GPT-4 in few-shot scenarios; GPT-4 demonstrates tangible gains when provided with few-shot prompts, whereas DeepSeekMath~displays consistent results across zero-shot and few-shot tests. Moving forward, our efforts will focus on enhancing the data selection process to create a superior pre-training corpus and, additionally, to pursue the research directions identified in Section~\ref{sec:effec_rl} for optimizing reinforcement learning strategies for large language models.

\end{document}